\section{Perturbative Matching in Minkowski Space and Non-perturbative Renormalization in Euclidean Space}

In a recent paper~\cite{Carlson:2017gpk}, the validity about the factorization of
the quasi-distribution into parton distribution as in Eq.~(\ref{eq:fact}) has been questioned.
In the original one-loop factorization
proof in Ref.~\cite{Xiong:2013bka}, the quasi and physical parton distributions were
calculated using the Lehmann-Symanzik-Zimmermann reduction formula in Minkowskian field theory~\cite{Sterman:1994ce} where all the internal loop integrals are Minkowskian.
In this form, all the infrared (IR) properties are very clear. On the other hand, because the quasi-distribution
has no time-dependence, it can be calculated from a Euclidean space state that evolves as
\begin{equation}
|p\rangle = e^{-\beta H}|p)
\end{equation}
in Monte Carlo simulations. Here $|p\rm{)}$ is a state with momentum $p$ and correct quantum number but not an eigenstate, whereas $H$ is the QCD Hamiltonian, $\beta$ is an evolution parameter (``Euclidean time''). In lattice QCD, nucleon matrix elements are obtained from the source-operator-sink correlation where the source, operator, and sink are separated by imaginary time. For large separation, this three-point correlation function is dominated by the contribution from the lowest physical state~\cite{Rothe:1992nt}.
Thus, the matrix elements are effectively calculated in an Euclidean-like field theory. However,
the equivalence of the Euclidean calculations with the original Minkowskian field theory is guaranteed
through the derivation of the formalism. This equivalence has been the basis for all the
Euclidean lattice calculations in the past.

It was shown, however, that the Feynman integrals in lattice perturbation theory
with external Minkowskian momenta do not have collinear divergences~\cite{Carlson:2017gpk}. Therefore, it appears that the quasi-distribtutions from lattice QCD do not have proper IR properties.
In a subsequent paper~\cite{Briceno:2017cpo}, it was explicitly shown that the lattice
calculation procedure leads to the Minkowskian theory matrix element with correct collinear divergences.
The question then becomes what the correct procedure is to recover the physical
result (or collinear divergences) with Minkowskian external momenta from the
Euclidean perturbation theory. The answer is that one cannot directly insert
the Minskowskian external momentum into the Euclidean integrals, as this procedure will
change the pole structures. The correct procedure is to finish the Euclidean
integral with Euclidean external momenta and then analytically continue the result to Minkowski spacetime.
To see that the nucleon matrix elements of the quasi-distribution calculated from lattice QCD does capture the correct IR physics in perturbation theory, one has to perform an analytical continuation of the
matrix elements after the Euclidean loop integrals are done, not before.

In the following, we provide an example to illustrate the above point: We calculate a one-loop integral appearing in the Euclidean quasi-distribution and show that the collinear divergence becomes manifest after an analytic continuation from $p_E^2$ to $p_M^2$ with $p_M^\mu$ and $p_E^\mu$ being the Minkowskian and Euclidean momenta. Actually the result is identical to its Minkowskian counterpart.


We first start with the Minkowskian integral, where $p_M^0 = \sqrt{p_z^2+m^2}$,
\beq \label{eq:mloop}
I_M = \int {d^4k_M\over (2\pi)^4}{(2\pi)p^z\delta(k^z_M-xp^z)\over (k_M^2-m^2+i\epsilon) \left[(p_M-k_M)^2+i\epsilon\right]}\ ,
\eeq
We focus on the physical region $0<x<1$ where the collinear divergence could exist.
For $I_M$, we first integrate over $k_M^0$ with the residue theorem, and then over the transverse components, to get the result
\begin{align} 	\label{eq: im}
I_M =& {i\over 8\pi}{1\over \sqrt{1+\rho_M}}\left[\ln{1+\sqrt{1+\rho_M}\over \rho_M}\right.\nonumber\\
&\left. +\ln{\sqrt{(1+\rho_M)(x^2+\rho_M)}+x+\rho_M\over 1-x}\right]\ ,
\end{align}
where
\beq
\rho_M = {m^2\over (p^z)^2}<1\ .
\eeq


For the Euclidean space integral, we first keep $p^\tau_E$ real,
\beq \label{eq:eloop}
I_E = \int {d^4k_E\over (2\pi)^4}{(2\pi)p^z\delta(k^z_E-xp^z)\over (k_E^2+m^2) (p_E-k_E)^2}  \ .
\eeq
Integrating over the poles, we find
\begin{align}	\label{eq: ie}
I_E =& {1\over 8\pi\sqrt{\rho_E-1-i\ep}} \left[\arctan{\rho_E+\rho_M-2(1-x)\over2 \sqrt{\rho_E-1-i\ep}(1-x)} \right.\nonumber\\
&\left.+ \arctan {\rho_E-\rho_M-2x\over2 \sqrt{(\rho_E-1-i\ep)(x^2+\rho_M)}}\right] \ ,
\end{align}
where
\beq \label{eq: rho}
\rho_E={(p^\tau)^2+(p^z)^2\over (p^z)^2}>1\ .
\eeq
We keep $-i\ep$ in the imaginary part of $\rho_E$ so that we can analytically continue it to $\rho_M$ through the lower half complex plane by making the replacement $p^\tau \rightarrow -ip^0$,
\beq \label{eq: cont}
\sqrt{\rho_E-1-i\ep} \to -i\sqrt{1+\rho_M} \ .
\eeq
We find that
\begin{align} \label{eq: iecont}
I'_E =&-{1\over8\pi}{1\over \sqrt{1+\rho_M}} \left[ {1\over2}\text{arctanh}\left(- {1\over\sqrt{1+\rho_M}}\right)\right.\nonumber\\
&\left. + {1\over2}\text{arctanh}\left(-{x+\rho\over\sqrt{(1+\rho_M)(x^2+\rho_M)}}\right)\right]\nonumber\\
=& -i I_M \ ,
\end{align}
which is exactly the Minkowskian matrix element with a factor of $i$ that is to be canceled in $dk^0_M=idk^4_E$.
If we take the massless limit $m\to0$, then both $I_M$ and $I'_E$ give
\beq \label{eq: oslimit}
I_M = {i\over 8\pi} \left(\ln {(p^z)^2\over m^2} + \ln {4x\over1-x} \right) \ ,
\eeq
with collinear divergence correctly characterized by $\ln m^2$.

Despite the analytical continuation to recover the correct collinear divergence, we find that the Euclidean matrix elements with Euclidean external momenta are still useful for renormalization. The quasi-distribution calculated in lattice QCD requires a non-perturbative renormalization, which can be done, e.g., in the RI/MOM scheme used in Refs.~\cite{Alexandrou:2017huk,Chen:2017mzz,Stewart:2017tvs}. To compute the renormalization factors in such a scheme, one needs to compute the vertex functions with external quarks and gluons that have deep-Euclidean momenta. After renormalization, the quasi-distribution is expected to have a well-defined continuum limit. It can then be matched perturbatively to the normal distributions in the $\overline{\rm MS}$ scheme, where the matching factor can be computed in Minkowski space.
